\chapter{Diccionario de Clases}
\label{cap:diccionario}

Se documenta cada clase del diagrama de diseño indicando su responsabilidad, sus atributos y sus
operaciones principales. La notación de visibilidad sigue el estándar UML: \texttt{-} privado,
\texttt{+} público, \texttt{\#} protegido.

En cada tabla, el bloque \textbf{Atributos} indica nombre, tipo y descripción, y el bloque
\textbf{Operaciones} presenta la firma completa y su descripción.

\section{Paquete «Seguridad y maestros»}
\label{sec:pkg_maestros}

\subsection{Clase \clase{Deposito}}

\textbf{Responsabilidad:} representa al terminal extraportuario, depósito o bodega que opera el
sistema. Es la raíz de la configuración: todos los maestros y operaciones le pertenecen.

{\estilotabla
\begin{longtable}{L{0.26\textwidth} L{0.12\textwidth} L{0.54\textwidth}}
\caption{Clase \clase{Deposito}: atributos y operaciones}\label{tab:cls_deposito}\\
\toprule
\multicolumn{3}{l}{\cellcolor{filacabecera}\textbf{Atributos}} \\
\midrule
\textbf{Nombre} & \textbf{Tipo} & \textbf{Descripción} \\
\midrule
\endfirsthead
\toprule
\textbf{Nombre} & \textbf{Tipo} & \textbf{Descripción} \\
\midrule
\endhead
\clase{idDeposito} & int & Identificador único del depósito. \\
\clase{rut} & String & Identificación tributaria de la organización. \\
\clase{razonSocial} & String & Nombre legal del terminal o bodega. \\
\clase{direccion} & String & Ubicación física del recinto. \\
\clase{zonaHoraria} & String & Zona horaria de operación. Relevante para la
internacionalización. \\
\midrule
\multicolumn{3}{l}{\cellcolor{filacabecera}\textbf{Operaciones}} \\
\midrule
\multicolumn{2}{L{0.40\textwidth}}{\clase{+ obtenerParametros() : Parametro[]}} &
Recupera la parametrización vigente del depósito. \\
\bottomrule
\end{longtable}
}

\subsection{Clase \clase{Usuario}}

\textbf{Responsabilidad:} representa a una persona con acceso al sistema, en web o móvil.
Concentra la autenticación y la resolución de permisos.

{\estilotabla
\begin{longtable}{L{0.26\textwidth} L{0.12\textwidth} L{0.54\textwidth}}
\caption{Clase \clase{Usuario}: atributos y operaciones}\label{tab:cls_usuario}\\
\toprule
\multicolumn{3}{l}{\cellcolor{filacabecera}\textbf{Atributos}} \\
\midrule
\textbf{Nombre} & \textbf{Tipo} & \textbf{Descripción} \\
\midrule
\endfirsthead
\toprule
\textbf{Nombre} & \textbf{Tipo} & \textbf{Descripción} \\
\midrule
\endhead
\clase{idUsuario} & int & Identificador único. \\
\clase{nombre} & String & Nombre completo del usuario. \\
\clase{email} & String & Correo electrónico; actúa como identificador de acceso. \\
\clase{passwordHash} & String & Resumen criptográfico de la contraseña. Nunca se almacena en
claro. \\
\clase{activo} & boolean & Indica si la cuenta está habilitada. \\
\clase{ultimoAcceso} & DateTime & Fecha y hora del último ingreso, para la bitácora. \\
\midrule
\multicolumn{3}{l}{\cellcolor{filacabecera}\textbf{Operaciones}} \\
\midrule
\multicolumn{2}{L{0.40\textwidth}}{\clase{+ autenticar(pass : String) : boolean}} &
Valida las credenciales y establece la sesión (CU-22). \\
\multicolumn{2}{L{0.40\textwidth}}{\clase{+ tienePermiso(codigo : String) : boolean}} &
Resuelve si el usuario puede ejecutar una acción, recorriendo sus roles. \\
\multicolumn{2}{L{0.40\textwidth}}{\clase{+ cambiarPassword(nueva : String) : void}} &
Actualiza la contraseña aplicando la función de resumen. \\
\bottomrule
\end{longtable}
}

\subsection{Clase \clase{Rol}}

\textbf{Responsabilidad:} agrupa un conjunto de permisos bajo un nombre, para asignarlos en bloque.
Permite definir perfiles a medida (administrador, jefe de faena, jefe tarjador, supervisor).

{\estilotabla
\begin{longtable}{L{0.26\textwidth} L{0.12\textwidth} L{0.54\textwidth}}
\caption{Clase \clase{Rol}: atributos y operaciones}\label{tab:cls_rol}\\
\toprule
\multicolumn{3}{l}{\cellcolor{filacabecera}\textbf{Atributos}} \\
\midrule
\textbf{Nombre} & \textbf{Tipo} & \textbf{Descripción} \\
\midrule
\endfirsthead
\toprule
\textbf{Nombre} & \textbf{Tipo} & \textbf{Descripción} \\
\midrule
\endhead
\clase{idRol} & int & Identificador único. \\
\clase{nombre} & String & Denominación del perfil. \\
\clase{descripcion} & String & Alcance del rol. \\
\midrule
\multicolumn{3}{l}{\cellcolor{filacabecera}\textbf{Operaciones}} \\
\midrule
\multicolumn{2}{L{0.40\textwidth}}{\clase{+ agregarPermiso(p : Permiso) : void}} &
Incorpora un permiso al rol. \\
\bottomrule
\end{longtable}
}

\subsection{Clase \clase{Permiso}}

\textbf{Responsabilidad:} representa la autorización atómica para ejecutar una funcionalidad
concreta.

{\estilotabla
\begin{longtable}{L{0.26\textwidth} L{0.12\textwidth} L{0.54\textwidth}}
\caption{Clase \clase{Permiso}: atributos}\label{tab:cls_permiso}\\
\toprule
\textbf{Nombre} & \textbf{Tipo} & \textbf{Descripción} \\
\midrule
\endfirsthead
\toprule
\textbf{Nombre} & \textbf{Tipo} & \textbf{Descripción} \\
\midrule
\endhead
\clase{idPermiso} & int & Identificador único. \\
\clase{codigo} & String & Clave interna de la funcionalidad protegida (p. ej.
\estado{REPORTE\_APROBAR}). \\
\clase{descripcion} & String & Explicación legible del permiso. \\
\bottomrule
\end{longtable}
}

\subsection{Clase \clase{Cliente}}

\textbf{Responsabilidad:} representa a la empresa importadora consignataria de la carga. Es el
destinatario del reporte.

{\estilotabla
\begin{longtable}{L{0.26\textwidth} L{0.12\textwidth} L{0.54\textwidth}}
\caption{Clase \clase{Cliente}: atributos y operaciones}\label{tab:cls_cliente}\\
\toprule
\multicolumn{3}{l}{\cellcolor{filacabecera}\textbf{Atributos}} \\
\midrule
\textbf{Nombre} & \textbf{Tipo} & \textbf{Descripción} \\
\midrule
\endfirsthead
\toprule
\textbf{Nombre} & \textbf{Tipo} & \textbf{Descripción} \\
\midrule
\endhead
\clase{idCliente} & int & Identificador único. \\
\clase{rut} & String & Identificación tributaria del importador. \\
\clase{razonSocial} & String & Nombre legal de la empresa. \\
\clase{emailContacto} & String & Correo principal de contacto. \\
\clase{telefono} & String & Teléfono de contacto. \\
\clase{activo} & boolean & Indica si el cliente está vigente. \\
\midrule
\multicolumn{3}{l}{\cellcolor{filacabecera}\textbf{Operaciones}} \\
\midrule
\multicolumn{2}{L{0.40\textwidth}}{\clase{+ obtenerDestinatarios() : String[]}} &
Devuelve las direcciones de correo a las que debe enviarse el reporte (CU-20). \\
\bottomrule
\end{longtable}
}

\subsection{Clase \clase{Trabajador}}

\textbf{Responsabilidad:} representa al personal operativo del depósito que conforma las
cuadrillas.

{\estilotabla
\begin{longtable}{L{0.26\textwidth} L{0.12\textwidth} L{0.54\textwidth}}
\caption{Clase \clase{Trabajador}: atributos y operaciones}\label{tab:cls_trabajador}\\
\toprule
\multicolumn{3}{l}{\cellcolor{filacabecera}\textbf{Atributos}} \\
\midrule
\textbf{Nombre} & \textbf{Tipo} & \textbf{Descripción} \\
\midrule
\endfirsthead
\toprule
\textbf{Nombre} & \textbf{Tipo} & \textbf{Descripción} \\
\midrule
\endhead
\clase{idTrabajador} & int & Identificador único. \\
\clase{rut} & String & Identificación del trabajador. \\
\clase{nombre} & String & Nombre completo. \\
\clase{cargo} & String & Función que desempeña (p. ej. jefe tarjador, estibador). \\
\clase{habilitado} & boolean & Indica si puede ser asignado a faenas. \\
\midrule
\multicolumn{3}{l}{\cellcolor{filacabecera}\textbf{Operaciones}} \\
\midrule
\multicolumn{2}{L{0.40\textwidth}}{\clase{+ estaDisponible(f : Date, h : Time) : boolean}} &
Verifica que no tenga otra faena asignada en la franja indicada (CU-09). \\
\bottomrule
\end{longtable}
}

\subsection{Clase \clase{Cuadrilla}}

\textbf{Responsabilidad:} agrupa a los cuatro o cinco trabajadores que ejecutan físicamente una
faena de desconsolidado, uno de los cuales opera la aplicación móvil.

{\estilotabla
\begin{longtable}{L{0.26\textwidth} L{0.12\textwidth} L{0.54\textwidth}}
\caption{Clase \clase{Cuadrilla}: atributos y operaciones}\label{tab:cls_cuadrilla}\\
\toprule
\multicolumn{3}{l}{\cellcolor{filacabecera}\textbf{Atributos}} \\
\midrule
\textbf{Nombre} & \textbf{Tipo} & \textbf{Descripción} \\
\midrule
\endfirsthead
\toprule
\textbf{Nombre} & \textbf{Tipo} & \textbf{Descripción} \\
\midrule
\endhead
\clase{idCuadrilla} & int & Identificador único. \\
\clase{nombre} & String & Denominación de la cuadrilla. \\
\clase{turno} & String & Turno de trabajo asociado. \\
\midrule
\multicolumn{3}{l}{\cellcolor{filacabecera}\textbf{Operaciones}} \\
\midrule
\multicolumn{2}{L{0.40\textwidth}}{\clase{+ agregarTrabajador(t : Trabajador) : void}} &
Incorpora un integrante a la cuadrilla. \\
\multicolumn{2}{L{0.40\textwidth}}{\clase{+ obtenerJefeTarjador() : Trabajador}} &
Devuelve al integrante designado como responsable de la aplicación móvil. \\
\multicolumn{2}{L{0.40\textwidth}}{\clase{+ estaDisponible(f : Date, h : Time) : boolean}} &
Verifica la disponibilidad del equipo completo para la franja indicada. \\
\bottomrule
\end{longtable}
}

\subsection{Clase \clase{TipoBulto}}

\textbf{Responsabilidad:} parametriza las unidades con que se registra la carga (caja, pallet,
atado, bobina), permitiendo adaptar el sistema a la operación de cada depósito.

{\estilotabla
\begin{longtable}{L{0.26\textwidth} L{0.12\textwidth} L{0.54\textwidth}}
\caption{Clase \clase{TipoBulto}: atributos}\label{tab:cls_tipobulto}\\
\toprule
\textbf{Nombre} & \textbf{Tipo} & \textbf{Descripción} \\
\midrule
\endfirsthead
\toprule
\textbf{Nombre} & \textbf{Tipo} & \textbf{Descripción} \\
\midrule
\endhead
\clase{idTipoBulto} & int & Identificador único. \\
\clase{nombre} & String & Denominación del tipo de bulto. \\
\clase{descripcion} & String & Detalle del tipo. \\
\clase{unidadMedida} & String & Unidad en que se cuantifica. \\
\bottomrule
\end{longtable}
}

\section{Paquete «Información de carga»}
\label{sec:pkg_carga}

\subsection{Clase \clase{Manifiesto}}

\textbf{Responsabilidad:} representa el documento oficial que declara la totalidad de la carga que
transporta una nave. Es la fuente de verdad de referencia del sistema.

{\estilotabla
\begin{longtable}{L{0.26\textwidth} L{0.12\textwidth} L{0.54\textwidth}}
\caption{Clase \clase{Manifiesto}: atributos y operaciones}\label{tab:cls_manifiesto}\\
\toprule
\multicolumn{3}{l}{\cellcolor{filacabecera}\textbf{Atributos}} \\
\midrule
\textbf{Nombre} & \textbf{Tipo} & \textbf{Descripción} \\
\midrule
\endfirsthead
\toprule
\textbf{Nombre} & \textbf{Tipo} & \textbf{Descripción} \\
\midrule
\endhead
\clase{idManifiesto} & int & Identificador único. \\
\clase{numeroManifiesto} & String & Folio oficial asignado por Aduanas. \\
\clase{nave} & String & Nombre de la nave que transporta la carga. \\
\clase{viaje} & String & Identificador del viaje. \\
\clase{fechaRecepcion} & DateTime & Momento en que el sistema descargó el documento. \\
\clase{origenDatos} & OrigenDatos & Vía por la que ingresó la información (XML, API o tarja
ciega). \\
\midrule
\multicolumn{3}{l}{\cellcolor{filacabecera}\textbf{Operaciones}} \\
\midrule
\multicolumn{2}{L{0.40\textwidth}}{\clase{+ cargarDesdeXML(xml : String) : boolean}} &
Interpreta el documento XML y construye la jerarquía de contenedores y líneas de carga
(CU-05). \\
\multicolumn{2}{L{0.40\textwidth}}{\clase{+ validarEsquema() : boolean}} &
Verifica que el XML cumpla el estándar esperado antes de procesarlo. \\
\multicolumn{2}{L{0.40\textwidth}}{\clase{+ agregarContenedor(c : Contenedor) : void}} &
Asocia un contenedor al manifiesto. \\
\bottomrule
\end{longtable}
}

\subsection{Clase \clase{Contenedor}}

\textbf{Responsabilidad:} representa la unidad física de transporte que será desconsolidada.

{\estilotabla
\begin{longtable}{L{0.26\textwidth} L{0.12\textwidth} L{0.54\textwidth}}
\caption{Clase \clase{Contenedor}: atributos y operaciones}\label{tab:cls_contenedor}\\
\toprule
\multicolumn{3}{l}{\cellcolor{filacabecera}\textbf{Atributos}} \\
\midrule
\textbf{Nombre} & \textbf{Tipo} & \textbf{Descripción} \\
\midrule
\endfirsthead
\toprule
\textbf{Nombre} & \textbf{Tipo} & \textbf{Descripción} \\
\midrule
\endhead
\clase{idContenedor} & int & Identificador único interno. \\
\clase{sigla} & String & Prefijo del propietario del contenedor (estándar ISO). \\
\clase{numero} & String & Número identificador del contenedor. \\
\clase{tipo} & String & Tipo de contenedor (\este{dry}, \este{reefer}, \este{open top}). \\
\clase{tamano} & String & Dimensión (20', 40', 40'HC). \\
\clase{nroSelloDeclarado} & String & Número de sello de seguridad según el manifiesto. Se contrasta
en terreno. \\
\clase{fechaArribo} & Date & Fecha estimada o efectiva de llegada. \\
\midrule
\multicolumn{3}{l}{\cellcolor{filacabecera}\textbf{Operaciones}} \\
\midrule
\multicolumn{2}{L{0.40\textwidth}}{\clase{+ agregarLinea(l : LineaCarga) : void}} &
Asocia una línea de carga al contenedor. \\
\multicolumn{2}{L{0.40\textwidth}}{\clase{+ totalBultosDeclarados() : int}} &
Suma las cantidades declaradas de todas sus líneas. \\
\bottomrule
\end{longtable}
}

\subsection{Clase \clase{LineaCarga}}

\textbf{Responsabilidad:} representa el lote de mercancía consignado a un cliente específico dentro
de un contenedor. \textbf{Es la clase donde se materializa el contraste entre lo declarado y lo
recibido.}

{\estilotabla
\begin{longtable}{L{0.26\textwidth} L{0.12\textwidth} L{0.54\textwidth}}
\caption{Clase \clase{LineaCarga}: atributos y operaciones}\label{tab:cls_lineacarga}\\
\toprule
\multicolumn{3}{l}{\cellcolor{filacabecera}\textbf{Atributos}} \\
\midrule
\textbf{Nombre} & \textbf{Tipo} & \textbf{Descripción} \\
\midrule
\endfirsthead
\toprule
\textbf{Nombre} & \textbf{Tipo} & \textbf{Descripción} \\
\midrule
\endhead
\clase{idLinea} & int & Identificador único. \\
\clase{blConocimiento} & String & Número del conocimiento de embarque (\este{Bill of Lading}). \\
\clase{descripcion} & String & Descripción de la mercancía. \\
\clase{marcas} & String & Marcas y contramarcas impresas en los bultos. \\
\clase{cantidadDeclarada} & int & Número de bultos según el manifiesto. Nulo en modalidad tarja
ciega. \\
\clase{cantidadRecibida} & int & Número de bultos efectivamente contados en terreno. \\
\clase{pesoKg} & decimal & Peso declarado del lote. \\
\midrule
\multicolumn{3}{l}{\cellcolor{filacabecera}\textbf{Operaciones}} \\
\midrule
\multicolumn{2}{L{0.40\textwidth}}{\clase{+ contrastarCantidad(recibida : int) : int}} &
Registra la cantidad recibida y devuelve la diferencia respecto de la declarada. \\
\multicolumn{2}{L{0.40\textwidth}}{\clase{+ tieneDiferencia() : boolean}} &
Indica si existe faltante o sobrante, disparando la creación de una incidencia. \\
\bottomrule
\end{longtable}
}

\subsection{Enumeración \clase{OrigenDatos}}

\clase{MANIFIESTO\_XML} $\cdot$ \clase{INTEGRACION\_API} $\cdot$ \clase{TARJA\_CIEGA}: identifica la
vía por la que ingresó la información de la carga. Condiciona si el contraste de cantidades es
aplicable.

\section{Paquete «Operación en terreno»}
\label{sec:pkg_terreno}

\subsection{Clase \clase{Faena}}

\textbf{Responsabilidad: clase central del modelo.} Representa la operación programada de
desconsolidado de un contenedor: su planificación, ejecución y cierre. Actúa como agregado raíz del
expediente probatorio.

{\estilotabla
\begin{longtable}{L{0.26\textwidth} L{0.12\textwidth} L{0.54\textwidth}}
\caption{Clase \clase{Faena}: atributos y operaciones}\label{tab:cls_faena}\\
\toprule
\multicolumn{3}{l}{\cellcolor{filacabecera}\textbf{Atributos}} \\
\midrule
\textbf{Nombre} & \textbf{Tipo} & \textbf{Descripción} \\
\midrule
\endfirsthead
\toprule
\textbf{Nombre} & \textbf{Tipo} & \textbf{Descripción} \\
\midrule
\endhead
\clase{idFaena} & int & Identificador único. \\
\clase{folio} & String & Número de faena visible para el usuario y en el reporte. \\
\clase{fechaProgramada} & Date & Día en que se ejecutará el desconsolidado. \\
\clase{horaInicio} & Time & Inicio de la franja horaria asignada. \\
\clase{horaTermino} & Time & Término de la franja horaria asignada. \\
\clase{estado} & EstadoFaena & Situación actual de la faena. \\
\clase{modalidadIngreso} & OrigenDatos & Vía por la que se obtuvo la información de la carga. \\
\midrule
\multicolumn{3}{l}{\cellcolor{filacabecera}\textbf{Operaciones}} \\
\midrule
\multicolumn{2}{L{0.40\textwidth}}{\clase{+ programar(f : Date, hi : Time, ht : Time) : void}} &
Fija fecha y franja horaria (CU-08). \\
\multicolumn{2}{L{0.40\textwidth}}{\clase{+ asignarCuadrilla(c : Cuadrilla,
jt : Trabajador) : void}} &
Asocia el equipo de trabajo y designa al responsable de la aplicación móvil (CU-09). \\
\multicolumn{2}{L{0.40\textwidth}}{\clase{+ iniciar() : void}} &
Marca el comienzo efectivo de la faena en terreno. \\
\multicolumn{2}{L{0.40\textwidth}}{\clase{+ registrarEtapa(tipo : TipoEtapa) : EtapaFaena}} &
Abre una nueva etapa del proceso. \\
\multicolumn{2}{L{0.40\textwidth}}{\clase{+ cerrar() : void}} &
Valida la completitud de la evidencia, finaliza la faena y dispara la generación del reporte
(CU-14, CU-17). \\
\multicolumn{2}{L{0.40\textwidth}}{\clase{+ cambiarEstado(e : EstadoFaena) : void}} &
Transiciona el estado respetando las reglas de negocio. \\
\bottomrule
\end{longtable}
}

\subsection{Clase \clase{EtapaFaena}}

\textbf{Responsabilidad:} representa cada uno de los momentos críticos del proceso en los que el
sistema exige evidencia. Impone el orden del flujo en terreno.

{\estilotabla
\begin{longtable}{L{0.26\textwidth} L{0.12\textwidth} L{0.54\textwidth}}
\caption{Clase \clase{EtapaFaena}: atributos y operaciones}\label{tab:cls_etapafaena}\\
\toprule
\multicolumn{3}{l}{\cellcolor{filacabecera}\textbf{Atributos}} \\
\midrule
\textbf{Nombre} & \textbf{Tipo} & \textbf{Descripción} \\
\midrule
\endfirsthead
\toprule
\textbf{Nombre} & \textbf{Tipo} & \textbf{Descripción} \\
\midrule
\endhead
\clase{idEtapa} & int & Identificador único. \\
\clase{tipo} & TipoEtapa & Momento del proceso que representa. \\
\clase{fechaHoraInicio} & DateTime & Marca de tiempo de apertura de la etapa. \\
\clase{fechaHoraFin} & DateTime & Marca de tiempo de cierre de la etapa. \\
\clase{observacion} & String & Comentario libre del operario. \\
\midrule
\multicolumn{3}{l}{\cellcolor{filacabecera}\textbf{Operaciones}} \\
\midrule
\multicolumn{2}{L{0.40\textwidth}}{\clase{+ agregarEvidencia(f : Fotografia) : void}} &
Asocia una fotografía a la etapa. \\
\multicolumn{2}{L{0.40\textwidth}}{\clase{+ estaCompleta() : boolean}} &
Verifica que la etapa cuente con la evidencia obligatoria antes de permitir el avance. \\
\bottomrule
\end{longtable}
}

\subsection{Clase \clase{SelloSeguridad}}

\textbf{Responsabilidad:} registra la verificación del sello del contenedor, uno de los elementos
de mayor valor probatorio de toda la faena.

{\estilotabla
\begin{longtable}{L{0.26\textwidth} L{0.12\textwidth} L{0.54\textwidth}}
\caption{Clase \clase{SelloSeguridad}: atributos y operaciones}\label{tab:cls_sello}\\
\toprule
\multicolumn{3}{l}{\cellcolor{filacabecera}\textbf{Atributos}} \\
\midrule
\textbf{Nombre} & \textbf{Tipo} & \textbf{Descripción} \\
\midrule
\endfirsthead
\toprule
\textbf{Nombre} & \textbf{Tipo} & \textbf{Descripción} \\
\midrule
\endhead
\clase{idSello} & int & Identificador único. \\
\clase{numeroDigitado} & String & Número que el operario leyó y digitó en terreno. \\
\clase{estado} & EstadoSello & Resultado de la verificación. \\
\clase{fechaVerificacion} & DateTime & Momento de la verificación. \\
\midrule
\multicolumn{3}{l}{\cellcolor{filacabecera}\textbf{Operaciones}} \\
\midrule
\multicolumn{2}{L{0.40\textwidth}}{\clase{+ verificarContraDeclarado(dec : String) :
EstadoSello}} &
Compara el número digitado contra el declarado en el manifiesto y determina si el sello está
intacto, roto o no coincide. \\
\bottomrule
\end{longtable}
}

\subsection{Clase \clase{Fotografia}}

\textbf{Responsabilidad: unidad de evidencia del sistema.} Encapsula la imagen capturada en terreno
junto con los metadatos que le otorgan valor probatorio, y gestiona su ciclo de compresión y
subida.

{\estilotabla
\begin{longtable}{L{0.26\textwidth} L{0.12\textwidth} L{0.54\textwidth}}
\caption{Clase \clase{Fotografia}: atributos y operaciones}\label{tab:cls_fotografia}\\
\toprule
\multicolumn{3}{l}{\cellcolor{filacabecera}\textbf{Atributos}} \\
\midrule
\textbf{Nombre} & \textbf{Tipo} & \textbf{Descripción} \\
\midrule
\endfirsthead
\toprule
\textbf{Nombre} & \textbf{Tipo} & \textbf{Descripción} \\
\midrule
\endhead
\clase{idFoto} & int & Identificador único. \\
\clase{uuid} & String & Identificador universal generado en el dispositivo, previo a la subida. \\
\clase{nombreArchivo} & String & Nombre del archivo de imagen. \\
\clase{url} & String & Dirección en el repositorio en la nube, una vez subida. \\
\clase{fechaHoraCaptura} & DateTime & Momento exacto de la captura. Elemento probatorio clave. \\
\clase{latitud} & decimal & Latitud de la captura. \\
\clase{longitud} & decimal & Longitud de la captura. \\
\clase{tamanoBytes} & long & Peso del archivo tras la compresión. \\
\clase{hash} & String & Resumen criptográfico de la imagen, para verificar su integridad. \\
\clase{estado} & EstadoFoto & Situación de la fotografía en su ciclo de vida. \\
\clase{motivoRechazo} & String & Razón indicada por el Supervisor al rechazarla (CU-19). \\
\midrule
\multicolumn{3}{l}{\cellcolor{filacabecera}\textbf{Operaciones}} \\
\midrule
\multicolumn{2}{L{0.40\textwidth}}{\clase{+ comprimir() : void}} &
Reduce el peso de la imagen en el dispositivo manteniendo legibilidad suficiente como
evidencia. \\
\multicolumn{2}{L{0.40\textwidth}}{\clase{+ subir() : boolean}} &
Transfiere la imagen al repositorio en la nube. \\
\multicolumn{2}{L{0.40\textwidth}}{\clase{+ marcarSubida(url : String) : void}} &
Confirma la subida y registra la dirección definitiva. \\
\multicolumn{2}{L{0.40\textwidth}}{\clase{+ marcarRechazada(motivo : String) : void}} &
Invalida la fotografía y registra el motivo, solicitando recaptura. \\
\bottomrule
\end{longtable}
}

\subsection{Clase \clase{Incidencia}}

\textbf{Responsabilidad:} documenta toda anomalía detectada en la carga. Es la información que
activa los seguros y delimita las responsabilidades entre naviera, terminal e importador.

{\estilotabla
\begin{longtable}{L{0.26\textwidth} L{0.12\textwidth} L{0.54\textwidth}}
\caption{Clase \clase{Incidencia}: atributos y operaciones}\label{tab:cls_incidencia}\\
\toprule
\multicolumn{3}{l}{\cellcolor{filacabecera}\textbf{Atributos}} \\
\midrule
\textbf{Nombre} & \textbf{Tipo} & \textbf{Descripción} \\
\midrule
\endfirsthead
\toprule
\textbf{Nombre} & \textbf{Tipo} & \textbf{Descripción} \\
\midrule
\endhead
\clase{idIncidencia} & int & Identificador único. \\
\clase{tipo} & TipoIncidencia & Naturaleza de la anomalía. \\
\clase{descripcion} & String & Relato del operario sobre lo observado. \\
\clase{cantidadAfectada} & int & Número de bultos comprometidos. \\
\clase{gravedad} & String & Clasificación de severidad del hecho. \\
\clase{fechaHora} & DateTime & Momento de la detección. \\
\midrule
\multicolumn{3}{l}{\cellcolor{filacabecera}\textbf{Operaciones}} \\
\midrule
\multicolumn{2}{L{0.40\textwidth}}{\clase{+ adjuntarFotografia(f : Fotografia) : void}} &
Asocia evidencia exclusiva a la incidencia. Obligatorio (RI-2). \\
\bottomrule
\end{longtable}
}

\subsection{Enumeraciones del paquete}

{\estilotabla
\begin{longtable}{L{0.18\textwidth} L{0.40\textwidth} L{0.36\textwidth}}

\caption{Enumeraciones del paquete «Operación en terreno»}
\label{tab:enum_terreno}
\\

\toprule
\textbf{Enumeración} & \textbf{Valores} & \textbf{Uso} \\
\midrule
\endfirsthead

\toprule
\textbf{Enumeración} & \textbf{Valores} & \textbf{Uso} \\
\midrule
\endhead

\clase{EstadoFaena} &
\estado{PROGRAMADA} $\cdot$ \estado{EN\_EJECUCION} $\cdot$ \estado{FINALIZADA} $\cdot$
\estado{ANULADA} &
Ciclo de vida de la operación. \\
\midrule

\clase{TipoEtapa} &
\estado{APERTURA} $\cdot$ \estado{SEPARACION} $\cdot$ \estado{INCIDENCIA} $\cdot$
\estado{CIERRE} &
Momentos críticos que exigen evidencia. \\
\midrule

\clase{EstadoSello} &
\estado{INTACTO} $\cdot$ \estado{ROTO} $\cdot$ \estado{NO\_COINCIDE} &
Resultado de la verificación del sello. \\
\midrule

\clase{EstadoFoto} &
\estado{PENDIENTE} $\cdot$ \estado{SUBIDA} $\cdot$ \estado{RECHAZADA} $\cdot$
\estado{REEMPLAZADA} &
Ciclo de vida de la evidencia, incluyendo el modo offline y la recaptura. \\
\midrule

\clase{TipoIncidencia} &
\estado{DANO} $\cdot$ \estado{FALTANTE} $\cdot$ \estado{SOBRANTE} $\cdot$
\estado{MAL\_EMBALADO} &
Clasificación de las anomalías. \\

\bottomrule

\end{longtable}
}

\section{Paquete «Reportes y publicación»}
\label{sec:pkg_reportes}

\subsection{Clase \clase{Reporte}}

\textbf{Responsabilidad:} consolida el expediente completo de una faena en un documento PDF
verificable y gestiona su ciclo de revisión, aprobación y publicación.

{\estilotabla
\begin{longtable}{L{0.26\textwidth} L{0.12\textwidth} L{0.54\textwidth}}
\caption{Clase \clase{Reporte}: atributos y operaciones}\label{tab:cls_reporte}\\
\toprule
\multicolumn{3}{l}{\cellcolor{filacabecera}\textbf{Atributos}} \\
\midrule
\textbf{Nombre} & \textbf{Tipo} & \textbf{Descripción} \\
\midrule
\endfirsthead
\toprule
\textbf{Nombre} & \textbf{Tipo} & \textbf{Descripción} \\
\midrule
\endhead
\clase{idReporte} & int & Identificador único. \\
\clase{folio} & String & Número visible del reporte. \\
\clase{version} & int & Número de versión, que se incrementa tras cada ciclo de recaptura. \\
\clase{fechaGeneracion} & DateTime & Momento de generación del borrador. \\
\clase{fechaAprobacion} & DateTime & Momento de la aprobación por el Supervisor. \\
\clase{estado} & EstadoReporte & Situación del reporte en su flujo. \\
\clase{urlPdf} & String & Ubicación del archivo PDF generado. \\
\midrule
\multicolumn{3}{l}{\cellcolor{filacabecera}\textbf{Operaciones}} \\
\midrule
\multicolumn{2}{L{0.40\textwidth}}{\clase{+ generarPDF() : Archivo}} &
Solicita la composición del documento a partir de los datos y fotografías de la faena (CU-17). \\
\multicolumn{2}{L{0.40\textwidth}}{\clase{+ versionar() : void}} &
Incrementa la versión preservando el historial tras una corrección. \\
\multicolumn{2}{L{0.40\textwidth}}{\clase{+ aprobar(s : Usuario, fh : DateTime) : void}} &
Registra quién aprobó el reporte y cuándo, habilitando su envío (CU-18). \\
\multicolumn{2}{L{0.40\textwidth}}{\clase{+ obtenerDestinatarios() : String[]}} &
Reúne los correos de todos los clientes involucrados en la faena. \\
\multicolumn{2}{L{0.40\textwidth}}{\clase{+ registrarEnvio(e : EnvioReporte) : void}} &
Deja constancia del despacho del documento. \\
\bottomrule
\end{longtable}
}

\subsection{Clase \clase{EnvioReporte}}

\textbf{Responsabilidad:} registra cada intento de despacho del reporte, exitoso o fallido, como
parte del expediente probatorio.

{\estilotabla
\begin{longtable}{L{0.26\textwidth} L{0.12\textwidth} L{0.54\textwidth}}
\caption{Clase \clase{EnvioReporte}: atributos y operaciones}\label{tab:cls_envioreporte}\\
\toprule
\multicolumn{3}{l}{\cellcolor{filacabecera}\textbf{Atributos}} \\
\midrule
\textbf{Nombre} & \textbf{Tipo} & \textbf{Descripción} \\
\midrule
\endfirsthead
\toprule
\textbf{Nombre} & \textbf{Tipo} & \textbf{Descripción} \\
\midrule
\endhead
\clase{idEnvio} & int & Identificador único. \\
\clase{fechaHoraEnvio} & DateTime & Momento del intento de envío. \\
\clase{destinatarios} & String & Direcciones a las que se despachó el documento. \\
\clase{estadoEnvio} & String & Resultado del intento (\estado{ENVIADO} / \estado{FALLIDO}). \\
\clase{mensajeError} & String & Detalle del fallo, cuando corresponde. \\
\clase{intentos} & int & Número de reintentos acumulados. \\
\midrule
\multicolumn{3}{l}{\cellcolor{filacabecera}\textbf{Operaciones}} \\
\midrule
\multicolumn{2}{L{0.40\textwidth}}{\clase{+ reintentar() : boolean}} &
Vuelve a despachar el documento tras un fallo. \\
\bottomrule
\end{longtable}
}

\subsection{Enumeración \clase{EstadoReporte}}

\estado{BORRADOR} $\cdot$ \estado{EN\_CORRECCION} $\cdot$ \estado{APROBADO} $\cdot$
\estado{ENVIADO}: refleja el flujo de control de calidad descrito en CU-18 y CU-19.

\section{Paquete «Controladores y servicios»}
\label{sec:pkg_controladores}

{\estilotabla
\begin{longtable}{L{0.17\textwidth} L{0.10\textwidth} L{0.37\textwidth} L{0.28\textwidth}}

\caption{Controladores, interfaces y servicios del diseño}
\label{tab:controladores}
\\

\toprule
\textbf{Clase} & \textbf{Estereo\-tipo} & \textbf{Responsabilidad} &
\textbf{Operaciones principales} \\
\midrule
\endfirsthead

\toprule
\textbf{Clase} & \textbf{Estereo\-tipo} & \textbf{Responsabilidad} &
\textbf{Operaciones principales} \\
\midrule
\endhead

\clase{Gestor\-Manifiesto} & \ster{control} &
Coordina la ingesta de información de carga por sus tres vías. &
\clase{importarManifiesto()}, \clase{recibirDesdeAPI()}, \clase{declararTarjaCiega()} \\
\midrule

\clase{Gestor\-Programacion} & \ster{control} &
Coordina la planificación de faenas y la asignación de recursos. &
\clase{listarContenedoresPendientes()}, \clase{programarFaena()}, \clase{asignarCuadrilla()},
\clase{validarDisponibilidadHorario()} \\
\midrule

\clase{GestorFaena} & \ster{control} &
Coordina la ejecución de la faena en terreno. Es el controlador con mayor tráfico de mensajes del
sistema. &
\clase{iniciarFaena()}, \clase{verificarSello()}, \clase{registrarLote()},
\clase{registrarIncidencia()}, \clase{registrarEvidencia()}, \clase{cerrarFaena()} \\
\midrule

\clase{Gestor\-Reportes} & \ster{control} &
Coordina la generación, revisión, aprobación y publicación del reporte. &
\clase{generarBorrador()}, \clase{rechazarFotografia()}, \clase{aprobarReporte()},
\clase{enviarReporte()} \\
\midrule

\clase{Sincronizador\-Offline} & \ster{control} &
Desacopla la captura de evidencia de su transmisión. Mantiene la cola de pendientes y gestiona los
reintentos. Es el único objeto que conoce el repositorio externo. &
\clase{encolarSubida()}, \clase{hayConectividad()}, \clase{sincronizarPendientes()},
\clase{persistirLocal()} \\
\midrule

\clase{IServicio\-Aduana} & \ster{interface} &
Abstrae la comunicación con la aduana, permitiendo la internacionalización sin modificar el resto
del sistema. &
\clase{descargarXML()} \\
\midrule

\clase{Servicio\-AduanaWS} & --- &
Implementación concreta del \este{web service} del Servicio Nacional de Aduanas de Chile. &
\clase{descargarXML()} \\
\midrule

\clase{IPasarela\-Correo} & \ster{interface} &
Abstrae el despacho de correo electrónico. &
\clase{enviar()} \\
\midrule

\clase{Servicio\-CorreoSMTP} & --- &
Implementación concreta sobre protocolo SMTP. &
\clase{enviar()} \\
\midrule

\clase{GeneradorPDF} & \ster{service} &
Compone el documento PDF a partir de datos y fotografías, destacando las incidencias. Corresponde
al rol que cumplía JasperReports en la versión original del sistema. &
\clase{armarReporte()}, \clase{destacarIncidencias()} \\
\midrule

\clase{Repositorio\-Fotos} & \ster{service} &
Almacenamiento externo de las imágenes capturadas en terreno. &
\clase{subir()}, \clase{obtenerURL()}, \clase{eliminar()} \\

\bottomrule

\end{longtable}
}
